\chapter{Espacios de funciones y convergencia uniforme}\label{cap:nuevo11}\label{chap:25}\label{chap:26}
\index{metrica uniforme@métrica uniforme}\index{convergencia!puntual}\index{convergencia!uniforme}

Una sucesión de funciones puede aproximar bien cada valor y, sin embargo, hacerlo mal en una región que cambia con el índice. La métrica uniforme mide el mayor error en todo el dominio y convierte las funciones en puntos de un nuevo espacio métrico. A partir de esa construcción estudiaremos la conservación de la continuidad, las series de funciones y el paso al límite en integrales y derivadas.

Usaremos completitud, continuidad, compacidad, propiedades del supremo y cálculo diferencial e integral de una variable. Cuando construyamos límites de funciones con valores en un espacio normado $V$, exigiremos que $V$ sea Banach. Para comparar funciones, su dominio no necesita inicialmente ninguna estructura: basta una distancia en el espacio de llegada. Salvo indicación contraria, los conjuntos considerados son no vacíos; si el dominio es vacío, existe una sola función con ese dominio y su espacio de funciones es el espacio métrico de un punto.

\section{Funciones acotadas y distancia del supremo}
\label{lectura:25}
\begin{definicion}
Sean $S$ un conjunto y $(Y,d)$ un espacio métrico. Una función $f:S\to Y$ es \emph{acotada} si $f(S)$ está contenido en alguna bola de radio finito. Denotamos por $B(S,Y)$ el conjunto de esas funciones. Para $f,g\in B(S,Y)$ definimos
\[
 d_\infty(f,g)=\sup_{s\in S}d(f(s),g(s)).
\]
\end{definicion}

Antes de verificar los axiomas de una métrica debemos comprobar que el supremo es finito. Si $f(S)\subset B(y_0,R)$ y $g(S)\subset B(y_1,T)$, la desigualdad triangular da
\[
 d(f(s),g(s))\le d(f(s),y_0)+d(y_0,y_1)+d(y_1,g(s))
 <R+d(y_0,y_1)+T.
\]
Por tanto, el conjunto cuyo supremo tomamos es no vacío, está acotado superiormente y está formado por números no negativos.

\begin{proposicion}\label{ch25:metrica}
La función $d_\infty$ es una métrica en $B(S,Y)$.
\end{proposicion}
\begin{proof}
La no negatividad se hereda de $d$. Si $d_\infty(f,g)=0$, entonces
$0\le d(f(s),g(s))\le 0$ para cada $s$, de modo que $f(s)=g(s)$ en todo $S$. Recíprocamente, si $f=g$, todas las distancias son cero y el supremo es cero. La simetría se obtiene de $d(f(s),g(s))=d(g(s),f(s))$.

Para la desigualdad triangular, fijamos $s\in S$ y escribimos
\[
 d(f(s),h(s))\le d(f(s),g(s))+d(g(s),h(s))
 \le d_\infty(f,g)+d_\infty(g,h).
\]
La expresión de la derecha no depende de $s$ y es una cota superior para todos los términos de la izquierda. La propiedad que caracteriza al supremo implica
$d_\infty(f,h)\le d_\infty(f,g)+d_\infty(g,h)$.
\end{proof}

\begin{observacion}[Un supremo no tiene que alcanzarse]
Si $f(s)=s$ y $g(s)=0$ en $S=(0,1)$, entonces $d_\infty(f,g)=1$, aunque $|f(s)-g(s)|<1$ para cada punto. Así, demostrar una desigualdad estricta punto por punto suele producir solamente una desigualdad no estricta después de tomar supremos. Para obtener finalmente un error menor que $\eps$ trabajaremos, por ejemplo, con $\eps/2$.
\end{observacion}

\begin{lema}\label{ch25:unionacotada}
Si una familia no vacía $\mathcal A\subset B(S,Y)$ está acotada para $d_\infty$, entonces $\bigcup_{f\in\mathcal A}f(S)$ está acotado en $Y$.
\end{lema}
\begin{proof}
Elegimos $f_0\in\mathcal A$. Como $\mathcal A$ está acotada, existe $M<\infty$ tal que $d_\infty(f,f_0)\le M$ para toda $f\in\mathcal A$. Como $f_0$ es una función acotada, existen $y_0\in Y$ y $R<\infty$ con $d(f_0(s),y_0)\le R$ para todo $s\in S$. En consecuencia,
\[
 d(f(s),y_0)\le d(f(s),f_0(s))+d(f_0(s),y_0)\le M+R
\]
para toda $f\in\mathcal A$ y todo $s\in S$. La misma bola, de centro $y_0$ y radio $M+R+1$, contiene todas las imágenes.
\end{proof}

Esta prueba proporciona directamente el control necesario. No hace falta construir una subsucesión con diámetros crecientes. Hay que distinguir dos afirmaciones: cada función de una familia puede ser acotada sin que la familia sea acotada en $d_\infty$. Las funciones constantes $f_n(s)=n$ ilustran esa diferencia.

\section{Por qué la completitud se transmite a las funciones}
\index{completitud!espacios de funciones}
\begin{teorema}\label{ch25:completitud}
Si $Y$ es completo, entonces $(B(S,Y),d_\infty)$ es completo.
\end{teorema}
\begin{estrategia}
La condición de Cauchy uniforme permite construir primero un límite en cada punto. Después debemos resolver dos problemas diferentes: demostrar que esa función límite es acotada y demostrar que la convergencia hacia ella es uniforme. Construir únicamente los límites puntuales no concluye la prueba.
\end{estrategia}
\begin{proof}
Sea $(f_n)$ una sucesión de Cauchy en $B(S,Y)$. Para $s\in S$ fijo tenemos
\[
 d(f_n(s),f_m(s))\le d_\infty(f_n,f_m).
\]
Por ello $(f_n(s))$ es de Cauchy en $Y$. La completitud de $Y$ proporciona un único punto $f(s)$ al cual converge esa sucesión. Al realizar esta construcción para cada $s$ obtenemos una función $f:S\to Y$.

Tomemos $N$ tal que $d_\infty(f_n,f_m)<1$ para $n,m\ge N$. Fijando $n=N$ y pasando al límite $m\to\infty$, la continuidad de la distancia da
\[
 d(f_N(s),f(s))\le 1 \qquad(s\in S).
\]
Si $f_N(S)\subset B(y_0,R)$, se sigue que $d(f(s),y_0)\le 1+R$ para todo $s$. Por tanto, $f\in B(S,Y)$.

Ahora sea $\eps>0$. Existe $N_\eps$ tal que $d_\infty(f_n,f_m)<\eps/2$ para $n,m\ge N_\eps$. Para $n\ge N_\eps$ y $s\in S$ fijos, al hacer $m\to\infty$ obtenemos
\[
 d(f_n(s),f(s))\le\eps/2.
\]
El índice $N_\eps$ fue elegido antes que $s$ y sirve para todos los puntos. Tomando supremos, $d_\infty(f_n,f)\le\eps/2<\eps$. Esto demuestra la convergencia en la métrica uniforme.
\end{proof}

\begin{ejemplo}[Por qué la hipótesis de completitud es necesaria]
Si $S=\{s_0\}$, la aplicación $f\mapsto f(s_0)$ es una isometría de $B(S,Y)$ sobre $Y$. Así, un espacio de llegada incompleto ya produce un espacio de funciones incompleto con un dominio de un solo punto. Más generalmente, para cualquier $S$ no vacío, las funciones constantes asociadas a una sucesión de Cauchy sin límite en $Y$ no pueden converger en $B(S,Y)$.
\end{ejemplo}

\section{Continuidad y límite uniforme}
\index{convergencia uniforme!continuidad del limite@continuidad del límite}
Si $X$ es un espacio métrico, escribimos $C_b(X,Y)$ para el conjunto de funciones continuas y acotadas de $X$ a $Y$. El subíndice $b$ recuerda la condición de acotación. Cuando $Y=\F$, usamos $C_b(X)$.

\begin{teorema}[Continuidad del límite uniforme]\label{ch25:cerradocb}\label{ch12:limite-continuo}
Sean $X,Y$ espacios métricos y $f_n:X\to Y$ funciones continuas. Si $f_n$ converge uniformemente a $f:X\to Y$, entonces $f$ es continua, sin necesidad de que las funciones sean acotadas. En particular, el conjunto $C_b(X,Y)$ es cerrado en $B(X,Y)$.
\end{teorema}
\begin{proof}
La convergencia uniforme significa que, dada una tolerancia positiva, todos los valores de $f_n$ están a esa distancia de los de $f$ cuando $n$ es suficientemente grande; el índice se elige antes que el punto. Para funciones acotadas equivale a convergencia en $d_\infty$. Fijemos $x\in X$ y $\eps>0$. Elegimos un índice $n$ tal que $d(f_n(z),f(z))<\eps/3$ para todo $z\in X$. Mantenemos fijo ese índice. La continuidad de $f_n$ en $x$ proporciona $\delta>0$ tal que
\[
 d_X(x,y)<\delta\quad\Longrightarrow\quad d(f_n(x),f_n(y))<\eps/3.
\]
Para esos puntos $y$, la desigualdad triangular da
\[
 \begin{split}
 d(f(x),f(y))&\le d(f(x),f_n(x))+d(f_n(x),f_n(y))+d(f_n(y),f(y))\\
 &<\eps/3+\eps/3+\eps/3=\eps.
 \end{split}
\]
Por tanto, $f$ es continua en $x$. Como $x$ era arbitrario, es continua en $X$. La elección uniforme del índice permite controlar a la vez el error en $x$ y en los puntos vecinos, sin exigir acotación global. Para la última afirmación, una sucesión de $C_b(X,Y)$ que converge en $B(X,Y)$ tiene un límite ya acotado por pertenecer a ese espacio ambiente y continuo por lo que acabamos de probar. La caracterización secuencial de cerrado permite concluir que $C_b(X,Y)$ es cerrado.
\end{proof}

El orden de las elecciones es decisivo: primero se fija una función $f_n$ que aproxima a $f$ en \emph{todos} los puntos y luego se usa la continuidad de esa función. No se necesita que el mismo $\delta$ funcione para todas las funciones de la sucesión.

\begin{corolario}
Si $Y$ es completo, entonces $C_b(X,Y)$ es completo con la métrica uniforme.
\end{corolario}
\begin{proof}
Una sucesión de Cauchy en $C_b(X,Y)$ converge en $B(X,Y)$ por el teorema~\ref{ch25:completitud}. Su límite pertenece a $C_b(X,Y)$ por el teorema~\ref{ch25:cerradocb}. Este es precisamente el significado de completitud del subespacio.
\end{proof}

\section{La norma uniforme y tres espacios que no deben confundirse}
\index{norma!uniforme}\index{funciones!de soporte compacto}\index{funciones!que se anulan en el infinito}
Sea $V$ un espacio normado sobre $\F$. Las operaciones puntuales
\[
 (f+g)(s)=f(s)+g(s),\qquad (\lambda f)(s)=\lambda f(s)
\]
convierten $B(S,V)$ en un espacio vectorial: las identidades vectoriales se verifican en cada punto; además, la suma y el producto por escalar conservan la acotación porque
$\norm{f(s)+g(s)}\le\norm{f(s)}+\norm{g(s)}$ y $\norm{\lambda f(s)}=|\lambda|\norm{f(s)}$.

\begin{proposicion}
La expresión $\norm f_\infty=\sup_{s\in S}\norm{f(s)}$ es una norma en $B(S,V)$ y su métrica asociada es $d_\infty$. Si $V$ es Banach, también lo son $B(S,V)$ y $C_b(X,V)$.
\end{proposicion}
\begin{proof}
La norma es no negativa y se anula exactamente cuando todos los valores $f(s)$ son cero. Para $\lambda\ne0$, multiplicar todas las cantidades de un conjunto por $|\lambda|$ multiplica su supremo por ese factor; así $\norm{\lambda f}_\infty=|\lambda|\norm f_\infty$. Para $\lambda=0$ la igualdad también se cumple. Finalmente,
\[
 \norm{f(s)+g(s)}\le\norm{f(s)}+\norm{g(s)}\le\norm f_\infty+\norm g_\infty;
\]
tomando supremos se obtiene la desigualdad triangular. La identidad $d_\infty(f,g)=\norm{f-g}_\infty$ se sigue de las definiciones. La completitud ya quedó demostrada. El conjunto $C_b(X,V)$ es un subespacio vectorial porque las operaciones de $V$ son continuas, y es cerrado por el teorema~\ref{ch25:cerradocb}.
\end{proof}

Para funciones escalares en un espacio métrico $X$ usaremos las notaciones
\[
 \begin{split}
 C_b(X)&=\{f:X\to\F:f\text{ es continua y acotada}\},\\
 C_c(X)&=\{f:X\to\F:f\text{ es continua y }\supp f\text{ es compacto}\},\\
 C_0(X)&=\{f:X\to\F:f\text{ es continua y }
 \{x:|f(x)|\ge\eps\}\text{ es compacto para cada }\eps>0\},
 \end{split}
\]
donde $\supp f=\cl{\{x:f(x)\ne0\}}$. Estos espacios satisfacen $C_c(X)\subset C_0(X)\subset C_b(X)$. Para justificar la segunda inclusión, el conjunto $K=\{|f|\ge1\}$ es compacto y $f$ es acotada en $K$ por continuidad; fuera de $K$, $|f|<1$. La primera inclusión se obtiene porque cada conjunto $\{|f|\ge\eps\}$ es cerrado y está contenido en el soporte compacto.

\begin{observacion}[Corrección de notación de la lectura original]
El espacio de funciones continuas de soporte compacto no es, en general, $C_b(X)$. En $X=\R$, la función constante $1$ pertenece a $C_b(\R)$ y no pertenece a $C_0(\R)$; la función $e^{-x^2}$ pertenece a $C_0(\R)$ y no pertenece a $C_c(\R)$. Si $X$ es compacto, en cambio, los tres espacios coinciden con $C(X)$.
\end{observacion}

\begin{proposicion}
Para un espacio métrico $X$, el espacio $C_0(X)$ es cerrado en $C_b(X)$. Por tanto, es Banach. Si, además, $X$ es localmente compacto, $C_c(X)$ es denso en $C_0(X)$ para la norma uniforme.
\end{proposicion}
\begin{proof}
Supongamos que $f_n\in C_0(X)$ y $\norm{f_n-f}_\infty\to0$. Dado $\eps>0$, elegimos $n$ con $\norm{f_n-f}_\infty<\eps/2$. Entonces
\[
 \{|f|\ge\eps\}\subset\{|f_n|\ge\eps/2\}.
\]
El conjunto de la izquierda es cerrado en $X$ y el de la derecha es compacto; por ello el primero es compacto. Esto demuestra $f\in C_0(X)$.

Para la densidad, fijemos $f\in C_0(X)$ y $\eps>0$. El conjunto $K=\{|f|\ge\eps/2\}$ es compacto. Si es vacío, la función cero aproxima a $f$ con error a lo sumo $\eps/2$. Si no es vacío, la compacidad local y una subcubierta finita proporcionan un abierto $U$ tal que $K\subset U$ y $\cl U$ es compacto. En efecto, cubrimos $K$ por finitos abiertos cuyas clausuras son compactas y tomamos su unión. Si $U=X$, podemos tomar $\chi=1$. En otro caso, los conjuntos cerrados disjuntos $K$ y $X\setminus U$ permiten definir
\[
 \chi(x)=\frac{\dist(x,X\setminus U)}{\dist(x,X\setminus U)+\dist(x,K)}.
\]
Las funciones distancia son continuas; el denominador no se anula, pues ello colocaría a $x$ en ambos conjuntos cerrados. Además, $0\le\chi\le1$, $\chi=1$ en $K$ y $\chi=0$ fuera de $U$. Por ello $h=\chi f$ es continua y su soporte está contenido en $\cl U$, luego es compacto. En $K$ tenemos $h=f$; fuera de $K$,
$|f-h|=|1-\chi|\,|f|<\eps/2$. Así $\norm{f-h}_\infty\le\eps/2<\eps$.
\end{proof}

\begin{observacion}[Verificaciones de la lectura 25]
La lectura 25 de la colección de partida no contiene una lista de ejercicios. Sus verificaciones están desarrolladas en las secciones anteriores: los axiomas de la métrica y de la norma, el control común de las imágenes, la completitud y la conservación de la continuidad. En particular, la prueba de completitud incluye explícitamente la convergencia uniforme hacia el límite construido.
\end{observacion}

\section{Convergencia puntual y uniforme: el orden de los cuantificadores}
\label{lectura:26}
La convergencia puntual responde una pregunta cada vez: después de fijar $x$, ¿se aproxima $f_n(x)$ a $f(x)$? La convergencia uniforme exige controlar simultáneamente todos los puntos con un único índice. Esa diferencia permite conservar continuidad, intercambiar integrales y límites y, con hipótesis adicionales, derivar límites.

La métrica anterior describe este control cuando las funciones son acotadas. Conviene formular ahora los cuantificadores de manera que también cubran funciones no acotadas y permitan precisar qué operaciones conservan los límites.

\begin{definicion}
Sean $f_n,f:S\to V$. Decimos que $f_n$ converge \emph{puntualmente} a $f$ si
\[
 \forall x\in S\ \forall\eps>0\ \exists N=N(x,\eps)\ \forall n\ge N:
 \norm{f_n(x)-f(x)}<\eps.
\]
Decimos que converge \emph{uniformemente}, y escribimos $f_n\rightrightarrows f$, si
\[
 \forall\eps>0\ \exists N=N(\eps)\ \forall n\ge N\ \forall x\in S:
 \norm{f_n(x)-f(x)}<\eps.
\]
\end{definicion}

Esta definición admite funciones no acotadas. Lo que acaba siendo acotado es el error $f_n-f$. Para funciones de $B(S,V)$ equivale a $\norm{f_n-f}_\infty\to0$. Exigir que \emph{todos} los errores, incluso los de los primeros términos, sean acotados añade una restricción innecesaria: modificar un número finito de términos no debe cambiar la convergencia. Adoptaremos la formulación mediante cuantificadores.

\begin{proposicion}
La convergencia uniforme implica convergencia puntual. El recíproco es falso.
\end{proposicion}
\begin{proof}
El índice que funciona para todos los puntos funciona en particular para cualquier punto fijado. Para refutar el recíproco tomamos $f_n(x)=x^n$ en $[0,1]$. Si $0\le x<1$, la sucesión geométrica $x^n$ tiende a cero; si $x=1$, permanece igual a uno. Por tanto, el límite puntual es
\[
 f(x)=\begin{cases}0,&0\le x<1,\\1,&x=1.\end{cases}
\]
Para cada $n$, $\sup_{0\le x<1}x^n=1$, de modo que $\norm{f_n-f}_\infty=1$. El error no tiende a cero. También podemos detectar el problema porque el límite es discontinuo, mientras que un límite uniforme de funciones continuas es continuo.
\end{proof}

\begin{observacion}[Corrección de la lectura original]
La observación 2 de la lectura 26 invierte esta implicación. La dirección válida es uniforme $\Rightarrow$ puntual; el ejemplo que la acompaña confirma precisamente que la otra dirección no vale.
\end{observacion}

\section{El teorema de Dini}
\index{Dini, teorema de}
\begin{teorema}[Dini]\label{ch26:dini}
Sean $X$ compacto y $(f_n)$ una sucesión de funciones continuas $X\to\R$ que converge puntualmente a una función continua $f$. Si la sucesión es creciente en todos los puntos, o decreciente en todos los puntos, entonces $f_n\rightrightarrows f$.
\end{teorema}
\begin{estrategia}
La convergencia puntual produce una cubierta por regiones donde el error es pequeño. La compacidad reduce esa cubierta a un número finito de regiones. La monotonía permite reemplazar sus distintos índices por uno solo.
\end{estrategia}
\begin{proof}
Supongamos primero que $f_n(x)\le f_{n+1}(x)$ para todo $x$. Al pasar al límite se obtiene $f_n(x)\le f(x)$. Así, las funciones continuas $g_n=f-f_n$ son no negativas, decrecientes y convergen puntualmente a cero.

Fijemos $\eps>0$ y definamos $U_n=\{x:g_n(x)<\eps/2\}$. Cada $U_n$ es abierto por continuidad y $U_n\subset U_{n+1}$ por monotonía. Para cada $x$, la convergencia $g_n(x)\to0$ asegura que $x\in U_n$ para algún $n$; por tanto, $(U_n)$ cubre $X$. Una subcubierta finita $U_{n_1},\ldots,U_{n_r}$ cubre $X$. Si $N=\max_j n_j$, el carácter creciente de la cubierta implica $X=U_N$, y también $X=U_n$ para $n\ge N$. Por ello
\[
 0\le f(x)-f_n(x)<\eps/2\quad(n\ge N,\ x\in X),
\]
lo que implica $\norm{f-f_n}_\infty\le\eps/2<\eps$. Si $(f_n)$ es decreciente, aplicamos la prueba anterior a $(-f_n)$, que es creciente y converge a $-f$.
\end{proof}

Las hipótesis tienen funciones diferentes. La sucesión $x^n$ en $[0,1]$ muestra que no basta con un límite puntual cualquiera: se necesita continuidad del límite. En $(0,1)$, la misma sucesión tiene límite continuo cero y es decreciente, pero no converge uniformemente; aquí falta compacidad. Sin monotonía podemos construir picos continuos de altura uno con soportes disjuntos que se acercan a cero; convergen puntualmente a cero, pero su norma uniforme permanece igual a uno. Por ejemplo, sirven los triángulos de altura uno centrados en $2^{-n}$ y de semiancho $2^{-n-2}$, extendidos por cero en $[0,1]$.

\section{Series de funciones y el criterio M}
\index{series de funciones}\index{Weierstrass!criterio M}
\begin{definicion}
La serie $\sum_{n=1}^\infty u_n$ converge puntual o uniformemente cuando lo hace su sucesión de sumas parciales $s_N=\sum_{n=1}^N u_n$. Converge \emph{absolutamente en cada punto} si $\sum_n\norm{u_n(x)}<\infty$ para cada $x$. Se dice que converge \emph{normalmente} si $\sum_n\norm{u_n}_\infty<\infty$.
\end{definicion}

La convergencia normal es una condición de control uniforme de las normas. No debe confundirse con la convergencia absoluta punto por punto, que permite que los controles dependan de $x$.

\begin{teorema}[Criterio M de Weierstrass]\label{ch26:criteriom}\label{ch12:weierstrass}
Sean $V$ Banach y $u_n:S\to V$. Supongamos que existen $M_n\ge0$ tales que $\norm{u_n(x)}\le M_n$ para todo $x\in S$ y que $\sum_n M_n<\infty$. Entonces $\sum_nu_n$ converge uniformemente y absolutamente en cada punto. Si $s$ es su suma,
\[
 \norm{s-s_N}_\infty\le\sum_{n>N}M_n.
\]
\end{teorema}
\begin{proof}
Para $q>p$ y cualquier $x$,
\[
 \norm{s_q(x)-s_p(x)}
 \le\sum_{n=p+1}^q\norm{u_n(x)}\le\sum_{n=p+1}^qM_n.
\]
La última expresión tiende a cero uniformemente respecto de $q>p$ cuando $p\to\infty$, pues la serie numérica converge. Para cada $x$, $(s_N(x))$ es de Cauchy y, por completitud, tiene límite $s(x)\in V$. Fijando $p=N$ y haciendo $q\to\infty$ obtenemos
$\norm{s(x)-s_N(x)}\le\sum_{n>N}M_n$. Tomar el supremo prueba simultáneamente la convergencia uniforme y la estimación del error. Por otra parte, $\sum_n\norm{u_n(x)}\le\sum_nM_n<\infty$, lo que da la convergencia absoluta.
\end{proof}

\begin{observacion}[La completitud no es decorativa]
La lectura original formula este criterio para un espacio normado sin explicitar completitud. En el espacio $c_{00}$ de sucesiones de soporte finito, con norma $\ell^1$, la serie de vectores $\sum_n 2^{-n}e_n$ satisface $\sum_n\norm{2^{-n}e_n}=1$, pero su supuesto límite tendría infinitas coordenadas no nulas y no pertenecería a $c_{00}$. Un dominio $S$ de un solo punto convierte este ejemplo en una serie de funciones. Por eso formulamos el criterio en un espacio de Banach.
\end{observacion}

\section{Integración de límites uniformes}
\index{integracion@integración!limite uniforme@límite uniforme}
\begin{teorema}\label{ch26:integracion}
Si $f_n:[a,b]\to\F$ son Riemann-integrables y $f_n\rightrightarrows f$, entonces $f$ es Riemann-integrable y
\[
 \int_a^bf(x)\,dx=\lim_{n\to\infty}\int_a^bf_n(x)\,dx.
\]
\end{teorema}
\begin{proof}
El caso $a=b$ es inmediato, de modo que supongamos $a<b$. Comenzamos con funciones reales. Elegimos un índice $n_0$ tal que $\norm{f-f_{n_0}}_\infty<1$. Como una función Riemann-integrable es acotada, $f$ también lo es.

Para una partición $P:a=x_0<\cdots<x_r=b$, denotamos por $U(h,P)$ y $L(h,P)$ las sumas superior e inferior de Darboux. Si $\norm{f-f_n}_\infty\le\delta$, en cada subintervalo $I$ se cumple
\[
 \sup_I f\le\sup_I f_n+\delta,
 \qquad \inf_I f\ge\inf_I f_n-\delta.
\]
Al multiplicar las diferencias por las longitudes y sumar, resulta
\[
 U(f,P)-L(f,P)\le U(f_n,P)-L(f_n,P)+2\delta(b-a).
\]
Dado $\eta>0$, elegimos primero $n$ con $\norm{f-f_n}_\infty<\eta/[4(b-a)]$. Después, por integrabilidad de $f_n$, elegimos una partición $P$ con $U(f_n,P)-L(f_n,P)<\eta/2$. La desigualdad anterior proporciona $U(f,P)-L(f,P)<\eta$. El criterio de Darboux demuestra que $f$ es integrable.

Para funciones complejas, las partes real e imaginaria de $f_n$ son integrables y convergen uniformemente a las partes correspondientes de $f$, pues $|\operatorname{Re}z|,|\operatorname{Im}z|\le|z|$. El argumento real demuestra la integrabilidad de ambas componentes y, por definición, la de $f$.

Finalmente, para una función integrable $h$ real o compleja, cada suma de Riemann satisface
\[
 \left|\sum_j h(\xi_j)(x_j-x_{j-1})\right|
 \le\sum_j|h(\xi_j)|(x_j-x_{j-1})
 \le(b-a)\norm h_\infty.
\]
El paso al límite da $|\int_a^b h|\le(b-a)\norm h_\infty$. Aplicado a $h=f_n-f$, obtenemos
\[
 \left|\int_a^bf_n-\int_a^bf\right|
 \le(b-a)\norm{f_n-f}_\infty\longrightarrow0.
\]
\end{proof}

\begin{corolario}[Integración término a término]
Si cada $u_n$ es Riemann-integrable en $[a,b]$ y $\sum_nu_n$ converge uniformemente a $s$, entonces
$\int_a^bs=\sum_n\int_a^bu_n$.
\end{corolario}
\begin{proof}
Las sumas parciales son integrables por linealidad. Aplicamos el teorema~\ref{ch26:integracion} a esas sumas y usamos $\int_a^b\sum_{n=1}^Nu_n=\sum_{n=1}^N\int_a^bu_n$.
\end{proof}

\section{Derivación de límites: la hipótesis correcta}
\index{derivacion@derivación!limite uniforme@límite uniforme}
La convergencia uniforme de funciones no controla sus derivadas. Por ejemplo, $f_n(x)=\sin(nx)/n$ converge uniformemente a cero en $\R$, pero $f_n'(0)=1$ para todo $n$, mientras que la derivada del límite es cero. El resultado válido comienza con convergencia uniforme de las \emph{derivadas} y fija además una constante de integración mediante el valor en un punto.

\begin{lema}[Desigualdad del valor medio para funciones complejas]\label{ch26:tvmdesigualdad}
Si $h:[u,v]\to\C$ es continua y es derivable, respecto de la variable real, en $(u,v)$, y $|h'(t)|\le M$, entonces $|h(v)-h(u)|\le M(v-u)$.
\end{lema}
\begin{proof}
Si $h(v)=h(u)$, la desigualdad es inmediata. En otro caso elegimos $\lambda\in\C$ con $|\lambda|=1$ y $\lambda(h(v)-h(u))=|h(v)-h(u)|$. La función real $\phi(t)=\operatorname{Re}(\lambda h(t))$ satisface el teorema del valor medio real. Existe $\xi\in(u,v)$ tal que
\[
 |h(v)-h(u)|=\phi(v)-\phi(u)=(v-u)\phi'(\xi)
 \le(v-u)|h'(\xi)|\le(v-u)M.
\]
\end{proof}

En general no existe un punto $\xi$ que satisfaga la igualdad compleja $h(v)-h(u)=h'(\xi)(v-u)$. Por ejemplo, para $h(t)=e^{it}$ en $[0,2\pi]$ el lado izquierdo es cero y la derivada nunca se anula. La desigualdad anterior es el sustituto adecuado para la demostración.

\begin{teorema}\label{ch26:derivacion}
Sean $a<b$ finitos y $f_n:(a,b)\to\C$ derivables. Supongamos que para algún $x_0\in(a,b)$ converge $(f_n(x_0))$ y que $f_n'\rightrightarrows g$ en $(a,b)$. Entonces existe $f$ con $f_n\rightrightarrows f$ en $(a,b)$, la función $f$ es derivable y $f'=g$.
\end{teorema}
\begin{proof}
La convergencia uniforme de $f_n'$ implica que, para $n,m$ suficientemente grandes, $\sup_{(a,b)}|f_n'-f_m'|$ es finito y pequeño. Aplicando el lema anterior a $h=f_n-f_m$ entre $x_0$ y $x$, obtenemos
\[
 |f_n(x)-f_m(x)|
 \le |f_n(x_0)-f_m(x_0)|+|x-x_0|\sup_{(a,b)}|f_n'-f_m'|
 \le |f_n(x_0)-f_m(x_0)|+(b-a)\sup_{(a,b)}|f_n'-f_m'|.
\]
La última cota no depende de $x$ y tiende a cero cuando $n,m\to\infty$. El criterio de Cauchy uniforme, demostrado en el ejercicio~\ref{ej:lectura-26-1}, produce $f_n\rightrightarrows f$.

Para demostrar la derivabilidad fijamos $c\in(a,b)$ y definimos los cocientes prolongados
\[
 q_n(x)=\begin{cases}
 \dfrac{f_n(x)-f_n(c)}{x-c},&x\ne c,\\[4pt]
 f_n'(c),&x=c.
 \end{cases}
\]
La derivabilidad de $f_n$ implica exactamente que $q_n$ es continua en $c$. Para $x\ne c$, el lema aplicado a $f_n-f_m$ proporciona
\[
 |q_n(x)-q_m(x)|\le\sup_{(a,b)}|f_n'-f_m'|.
\]
La misma desigualdad vale para $x=c$ por la definición del cociente prolongado. Por tanto, $(q_n)$ es de Cauchy uniformemente y converge uniformemente a alguna función $q$.

El límite $q$ es continuo en $c$. Para verificarlo sin exigir continuidad de $q_n$ en otros puntos, dado $\eps>0$ elegimos $n$ con $\sup|q-q_n|<\eps/3$ y luego $\delta>0$ con $|q_n(x)-q_n(c)|<\eps/3$ si $|x-c|<\delta$. La desigualdad triangular da $|q(x)-q(c)|<\eps$.

Si $x\ne c$, la convergencia puntual de $f_n$ implica
$q(x)=(f(x)-f(c))/(x-c)$. En $x=c$ tenemos
$q(c)=\lim_n f_n'(c)=g(c)$. Por continuidad de $q$ en $c$,
\[
 \lim_{x\to c}\frac{f(x)-f(c)}{x-c}=q(c)=g(c).
\]
Así $f'(c)=g(c)$, y $c$ era arbitrario.
\end{proof}

No hemos supuesto que las derivadas sean continuas. Tampoco se obtiene una afirmación sobre valores arbitrariamente asignados en los extremos $a,b$: el teorema se refiere al intervalo abierto. La finitud de su longitud interviene en el factor $b-a$.

\section{Ejercicios sobre límites, productos y series uniformes}
\begin{ejercicio}[Lectura 26, ejercicio 1]\label{ej:lectura-26-1}\label{ch26:ej1}
Sea $V$ Banach y $f_n:S\to V$. Demuestre que $(f_n)$ converge uniformemente si y solo si para todo $\eps>0$ existe $N$ tal que $\norm{f_n(x)-f_m(x)}<\eps$ para todos $n,m\ge N$ y $x\in S$.
\end{ejercicio}


\begin{ejercicio}[Lectura 26, ejercicio 2]\label{ej:lectura-26-2}
Sean $(f_n)$ y $(g_n)$ sucesiones de funciones escalares sobre $S$, acotadas en $B(S,\F)$, que convergen uniformemente. Demuestre que $(f_ng_n)$ converge uniformemente. Dé un contraejemplo al suprimir el control de acotación.
\end{ejercicio}


\begin{ejercicio}[Lectura 26, ejercicio 3: enunciado corregido]\label{ej:lectura-26-3}
Sean $V,W$ Banach, $f_n\rightrightarrows f$ y $\norm{f_n}_\infty\le\alpha$. La lectura propone que la continuidad de $g$ en la bola de radio $\alpha$ basta para que $g\circ f_n\rightrightarrows g\circ f$. Examine la afirmación y pruebe una versión válida.
\end{ejercicio}


\begin{ejercicio}[Lectura 26, ejercicio 4]\label{ej:lectura-26-4}
Sean $V$ Banach y $f_n:S\to V$. Suponga que $\norm{f_{n+1}(x)-f_n(x)}\le M_n$ para todo $x$, donde $\sum_nM_n<\infty$. Demuestre que $(f_n)$ converge uniformemente.
\end{ejercicio}


\begin{ejercicio}[Lectura 26, ejercicio 5]\label{ej:lectura-26-5}
Estudie la relación entre convergencia absoluta y uniforme:
\begin{enumerate}[label=(\roman*)]
\item En el dominio $[0,1]$, para $x>0$ sea $u_n(x)=\sin(\pi/x)\mathbf1_{[1/(n+1),1/n]}(x)$ y defina $u_n(0)=0$. Analice la convergencia absoluta puntual de $\sum_nu_n$ y la convergencia uniforme.
\item Demuestre que $\sum_{n=1}^\infty(-1)^n(x^2+n)/n^2$ converge uniformemente en $[0,1]$, pero no absolutamente en ningún punto.
\end{enumerate}
\end{ejercicio}


\begin{ejercicio}[Lectura 26, ejercicio 6]\label{ej:lectura-26-6}
Si $f_n:X\to V$ son continuas y convergen uniformemente a $f$, demuestre que $f_n(x_n)\to f(x)$ siempre que $x_n\to x$. ¿Vale el recíproco?
\end{ejercicio}


\begin{ejercicio}[Lectura 26, ejercicio 7: Dirichlet uniforme]\label{ej:lectura-26-7}
Sean $V$ Banach, $u_n:S\to V$ y $v_n:S\to\R$. Suponga que, para algún $M<\infty$,
\[
 \sup_{x\in S}\left\|\sum_{k=1}^nu_k(x)\right\|\le M\quad(n\ge1),
\]
que $v_n\rightrightarrows0$ y que $v_{n+1}(x)\le v_n(x)$ para todo $x$. Demuestre que $\sum_nv_nu_n$ converge uniformemente.
\end{ejercicio}


\begin{ejercicio}[Lectura 26, ejercicio 8: condensación uniforme de Cauchy]\label{ej:lectura-26-8}
Sea $(u_n)$ una sucesión de funciones reales sobre $S$, decreciente respecto de $n$ en cada punto. Demuestre que $\sum_{n=1}^\infty u_n$ y $\sum_{k=1}^\infty 2^ku_{2^k}$ convergen uniformemente o dejan de hacerlo simultáneamente.
\end{ejercicio}


\section{El problema de Cauchy como ecuación entre funciones}\label{ampl:picard}
Hasta ahora hemos construido límites de funciones y comprobado qué propiedades conservan. Podemos emplear esos resultados para construir una función desconocida a partir de una ecuación. En un problema de Cauchy se prescribe una velocidad que depende del tiempo y del estado, junto con el estado inicial:
\[
 u'(t)=f(t,u(t)),\qquad u(t_0)=\xi.
\]
La idea es integrar antes de buscar la solución. Así la incógnita aparece como punto fijo de una transformación sobre un espacio completo de funciones, y el teorema~\ref{ch11:banach} del capítulo~\ref{cap:nuevo04} se vuelve aplicable.

Las dos demostraciones siguientes retoman las estrategias de Maximenko \cite{max-picard}. Aquí trabajamos en un intervalo compacto y permitimos valores en $\R^d$. No exigimos que $f$ sea acotada en todo el cilindro $I\times\R^d$: la continuidad basta para que su composición con cada curva continua sea acotada. Esta precisión permite incluir ecuaciones lineales como $u'=u$.
\index{Picard, teorema de}\index{problema de Cauchy}\index{norma!ponderada}\index{aproximaciones sucesivas}

\begin{teorema}[Picard en un intervalo compacto]\label{ampl:picard-teorema}
Sea $I=[a,b]$, $a<b$, y sea $f:I\times\R^d\to\R^d$ continua. Suponga que existe $L\ge0$ tal que, para todo $t\in I$ y todo $x,y\in\R^d$,
\[
 \norm{f(t,x)-f(t,y)}_\infty\le L\norm{x-y}_\infty.
\]
Dados $t_0\in I$ y $\xi\in\R^d$, existe una única solución $u\in C^1(I,\R^d)$ del problema de Cauchy. En los extremos del intervalo se entienden derivadas laterales. Desde cualquier $u_0\in C(I,\R^d)$, las aproximaciones
\[
 u_{n+1}(t)=\xi+\int_{t_0}^t f(s,u_n(s))\,ds
\]
convergen uniformemente a esa solución.
\end{teorema}

\begin{proof}[Preparación común a las dos demostraciones]
El espacio $X=C(I,\R^d)$ es Banach para
\[
 \norm u_\infty=\sup_{t\in I}\norm{u(t)}_\infty
\]
por los teoremas~\ref{ch25:completitud} y~\ref{ch25:cerradocb}; $\R^d$ es completo y toda curva continua sobre $I$ es acotada. La integral vectorial se define integrando cada coordenada. La desigualdad escalar para integrales y el máximo de las coordenadas dan
\[
 \left\|\int_c^d v(s)\,ds\right\|_\infty
 \le\int_{\min(c,d)}^{\max(c,d)}\norm{v(s)}_\infty\,ds.
\]
Para $u\in X$, la función $s\mapsto f(s,u(s))$ es continua. Por el teorema fundamental del cálculo,
\[
 (Su)(t)=\xi+\int_{t_0}^t f(s,u(s))\,ds
\]
pertenece a $C^1(I,\R^d)$ y, en particular, a $X$. Un punto fijo de $S$ es una solución del problema de Cauchy; recíprocamente, integrar la ecuación de cualquier solución muestra que es punto fijo. Es suficiente probar existencia, unicidad y convergencia de las iteraciones de $S$. Si $T=\max\{t_0-a,b-t_0\}$, la estimación directa sólo da
\[
 \norm{Su-Sv}_\infty\le LT\norm{u-v}_\infty.
\]
Cuando $LT\ge1$ esta cota no permite aplicar Banach directamente. No significa que la solución deje de existir; significa que necesitamos una estimación mejor.
\end{proof}

\begin{proof}[Primera demostración: una norma ponderada]
Elijamos $\beta>0$ con $\beta\ge L$ y definamos
\[
 \norm u_\beta=\sup_{t\in I}e^{-\beta|t-t_0|}\norm{u(t)}_\infty.
\]
Es una norma: la positividad del peso asegura que sólo se anula en la función cero, y homogeneidad y desigualdad triangular se verifican punto por punto antes de tomar el supremo. Además,
\begin{equation}\label{ampl:picard-equivalencia}
 e^{-\beta T}\norm u_\infty\le\norm u_\beta\le\norm u_\infty.
\end{equation}
Así una sucesión de Cauchy para la norma ponderada lo es para la uniforme, y su límite uniforme también es límite en la ponderada. Por tanto $X$ sigue siendo completo.

Para $r=|t-t_0|$, parametrizamos el segmento entre $t_0$ y $t$ por su distancia $s$ al punto inicial; esto evita cambiar signos cuando $t<t_0$. Obtenemos
\[
 \begin{split}
 e^{-\beta r}\norm{Su(t)-Sv(t)}_\infty
 &\le L e^{-\beta r}\norm{u-v}_\beta\int_0^r e^{\beta s}\,ds\\
 &=\frac L\beta(1-e^{-\beta r})\norm{u-v}_\beta.
 \end{split}
\]
Al tomar supremos, $S$ es una contracción con constante
\begin{equation}\label{ampl:picard-constante}
 q=\frac L\beta(1-e^{-\beta T})<1.
\end{equation}
Si $L=0$, $q=0$; si $L>0$, usamos $L/\beta\le1$ y $1-e^{-\beta T}<1$. Banach da un único punto fijo y convergencia de todas las iteraciones en $\norm{\cdot}_\beta$. La equivalencia~\eqref{ampl:picard-equivalencia} convierte esa convergencia en uniforme. La preparación anterior identifica el punto fijo con la solución requerida.
\end{proof}

\begin{proof}[Segunda demostración: una potencia contractiva]
Demostraremos por inducción que, para todo $n\ge0$,
\begin{equation}\label{ampl:picard-factorial}
 \norm{S^nu(t)-S^nv(t)}_\infty
 \le\frac{L^n|t-t_0|^n}{n!}\norm{u-v}_\infty.
\end{equation}
Para $n=0$, la afirmación es la definición de norma uniforme; interpretamos $L^0=1$. Supongamos cierta la cota para $n$. Al aplicar $S$ una vez más y parametrizar de nuevo el segmento desde $t_0$,
\[
 \norm{S^{n+1}u(t)-S^{n+1}v(t)}_\infty
 \le L\int_0^{|t-t_0|}\frac{L^ns^n}{n!}\norm{u-v}_\infty\,ds,
\]
que es la cota deseada con $n+1$. En consecuencia,
\[
 \norm{S^nu-S^nv}_\infty\le\frac{(LT)^n}{n!}\norm{u-v}_\infty.
\]
Si $L=0$, $S$ ya es constante. Si $L>0$, los números $(LT)^n/n!$ tienden a cero: a partir de un índice, el cociente entre dos consecutivos es $LT/(n+1)\le1/2$. Existe, pues, $N\ge1$ para el cual $S^N$ es contracción.

Sea $u$ su único punto fijo. Como $S$ conmuta con $S^N$,
$S^N(Su)=S(S^Nu)=Su$; así $Su$ es otro punto fijo de $S^N$ y, por unicidad, $Su=u$. Cualquier punto fijo de $S$ lo es de $S^N$, por lo que también es único. Esto desarrolla el mecanismo del ejercicio~\ref{ej:lectura-11-1}.

Finalmente, para cualquier $u_0$, Banach da $S^{kN}u_0\to u$. Para cada residuo $0\le j<N$, la aplicación $S^j$ es continua por~\eqref{ampl:picard-factorial}; luego
$S^{kN+j}u_0\to S^ju=u$. Hay sólo un número finito de residuos, de modo que las $N$ subsucesiones juntas forman una sucesión que converge a $u$. La convergencia es uniforme y concluye la segunda prueba.
\end{proof}

\begin{proposicion}[Un error que puede calcularse]\label{ampl:picard-error}
En el teorema anterior, inicie con $u_0(t)=\xi$ y ponga
$M=\max_{t\in I}\norm{f(t,\xi)}_\infty$. Para $n\ge0$,
\[
 \norm{u-u_n}_\infty
 \le MT\sum_{k=n}^{\infty}\frac{(LT)^k}{(k+1)!}
 \le MT e^{LT}\frac{(LT)^n}{(n+1)!}.
\]
Cuando $L=0$, las iteraciones son exactas desde $u_1$; para $n=0$ se interpreta $(LT)^0=1$.
\end{proposicion}
\begin{proof}
La primera diferencia satisface $\norm{u_1(t)-u_0(t)}_\infty\le M|t-t_0|$. El mismo argumento inductivo por integración da
\[
 \norm{u_{k+1}(t)-u_k(t)}_\infty
 \le\frac{ML^k|t-t_0|^{k+1}}{(k+1)!}.
\]
Sumar desde $k=n$ hasta $m-1$ y dejar $m\to\infty$ produce la primera cota, pues $u_m\to u$ uniformemente. Para la segunda, escribimos $k=n+j$ y usamos
$(n+j+1)!\ge(n+1)!j!$. La suma restante queda acotada por $\sum_{j\ge0}(LT)^j/j!=e^{LT}$. La cota muestra por qué las integraciones repetidas pueden mejorar mucho el control: aparece un factorial en el denominador.
\end{proof}

\section{Ejercicios sobre existencia, iteración y estabilidad}
Los cinco ejercicios de esta sección son desarrollos propios de esta edición a partir del método explicado en \cite{max-picard}; no son una lista reproducida de ese apunte.

\begin{ejercicio}[Ampliación: iteraciones que revelan la solución]\label{ej:ampl-picard-1}
Resuelva $u'=u+t$, $u(0)=1$, en $[0,2]$ mediante Picard. Partiendo de $u_0=1$, encuentre una fórmula para $u_n$ y una cota uniforme explícita del error. Explique por qué la estimación directa del operador no basta para aplicar Banach en la norma uniforme.
\end{ejercicio}


\begin{ejercicio}[Ampliación: una ecuación no lineal en un intervalo largo]\label{ej:ampl-picard-2}
Para $u'=\sin u+t$, $u(0)=0$, en $[0,2]$, verifique las hipótesis de Picard, calcule $u_1,u_2$ desde $u_0=0$ y obtenga cotas del error mediante la norma $\norm v_1=\sup_{0\le t\le2}e^{-t}|v(t)|$ y mediante el factorial.
\end{ejercicio}


\begin{ejercicio}[Ampliación: estabilidad frente a datos perturbados]\label{ej:ampl-picard-3}
Sean $u,v$ las soluciones en $I$ de $u'=f(t,u)$, $v'=g(t,v)$ con $u(t_0)=\xi$, $v(t_0)=\eta$. Suponga las hipótesis de Picard para ambas ecuaciones, que $f$ tiene constante Lipschitz $L$, y que $\norm{f(t,x)-g(t,x)}_\infty\le\delta$ para todo $t,x$. Demuestre, con $r=|t-t_0|$, que para $L>0$,
\[
 \norm{u(t)-v(t)}_\infty
 \le e^{Lr}\norm{\xi-\eta}_\infty
       +\delta\frac{e^{Lr}-1}{L}.
\]
Obtenga la versión para $L=0$ y deduzca dependencia continua respecto del dato inicial.
\end{ejercicio}


\begin{ejercicio}[Ampliación: versión local y explosión de una solución]\label{ej:ampl-picard-4}
Sea $f$ continua en $[t_0-a,t_0+a]\times\{x:\norm{x-\xi}_\infty\le b\}$, con $a,b>0$, acotada allí por $M$ y Lipschitz en $x$ con constante $L$. Si $0<h\le a$, $Mh\le b$ y $Lh<1$, pruebe existencia y unicidad en $[t_0-h,t_0+h]$ entre las soluciones que permanecen en esa bola. Aplique el resultado a $u'=u^2$, $u(0)=1$, y explique por qué no puede concluirse existencia en todo $[0,2]$.
\end{ejercicio}


\begin{ejercicio}[Ampliación: continuidad sin unicidad]\label{ej:ampl-picard-5}
Para el problema $u'=2\sqrt{|u|}$, $u(0)=0$, en $[0,1]$, pruebe que para cada $c\in[0,1]$,
\[
 u_c(t)=\begin{cases}0,&0\le t\le c,\\(t-c)^2,&c\le t\le1\end{cases}
\]
es una solución. Identifique exactamente qué hipótesis de unicidad falta, incluso si se restringe el estado a un intervalo compacto alrededor de cero.
\end{ejercicio}


\section{Cómo elegir una herramienta de convergencia}
La convergencia uniforme conserva la continuidad y permite pasar integrales al límite en intervalos compactos. Para derivar, se necesita controlar las derivadas y fijar un valor. El criterio de Cauchy construye el límite sin conocer su fórmula; Dini añade orden y compacidad a la convergencia puntual. El criterio M usa una cota sumable, mientras Dirichlet aprovecha cancelaciones. Estas diferencias determinan qué hipótesis deben comprobarse. Los capítulos~\ref{cap:nuevo12} y~\ref{cap:nuevo13} aplicarán este control a la compacidad, la extensión y la aproximación de funciones.
